\documentclass[11pt,twoside,a4paper]{article}
\title{Parameters and choices for AddiVortes}
\author{Lewis Maltby}
\usepackage{amsmath}
\usepackage[margin=0.8in]{geometry}
\begin{document}
\maketitle
\begin{center}
\begin{tabular}{ |c|c|c|c| }
\hline
\textbf{Choice/parameter} & \textbf{Symb} & \textbf{Default value} & \textbf{Search range} \\
\hline\hline
Metric to use & $d$ & Euclidean&\\
Number of tessellations & $m$ & 200 & [20, 500]\\
\hline\hline
Variance of noise & $\sigma$ & $\propto \frac{ \nu \lambda }{ \chi^2 _\nu }$\footnotemark[1]&\\[0.1cm]
\hline
Noise degrees of freedom hyperparam & $\nu$ & $3-10(6)$ & [1, 10]\\
Noise alignment hyperparam & $q$ & $0.7-0.99(0.85)$ & [0.6,0.999]\\
\hline\hline
Number of covariates in $T_j$ & $d_j$ & $d_j-1\sim Bin(p-1,\frac{\omega}{p})$& \\
Number of centers in $T_j$ & $b_j$ &$b_j-1\sim Po(\lambda_c)$& \\
Covariates included in $T_j$ &--- & Uniformly at random& \\
Coordinates of centers of $T_j$ & --- & $\mathcal{N}(0, \sigma^2_c)$\footnotemark[2]&\\
\hline
Number of covariates hyperparam & $\omega$ & $\min(3, p)$ & [1, 7] (depends on $p$) \\
Number of centers hyperparam & $\lambda_c$  & & [1, 50] \\
Center coordinates hyperparam & $\sigma_c$ & & [0.1, 3]\\
\hline
\hline
Mean response of cell $i$ under tess $j$ & $\mu_{ij}$ & $\mathcal{N}(\mu_\mu, \sigma_\mu^2)$&\\
\hline
Mean response center hyperparam & $\mu_\mu$ & 0\footnotemark[3]&\\
Mean response spread hyperparam & $\sigma_\mu$ & $\frac{0.5}{k\sqrt{m}}$\footnotemark[4]&\\
Mean response alignment hyperparam & $k$ & $3$&\\
\hline\hline
Prob of adding/removing a center & --- & 0.2&\\
Prob of adding/removing a covariate & --- & 0.2&\\
Prob of swapping a covariate & --- & 0.1&\\
Prob of changing position of a center & --- & 0.1&\\
\hline\hline
Number of MCMC iterations & --- & 1200 & [500, 100000] \\
Number of MCMC burn-ins & --- & 200 & [0.01,0.5] (as a proportion) \\
\hline
\end{tabular}
\end{center}
\footnotetext[1]{Two plausible choices for $\hat{\sigma}$ are considered. The first is the
“naive” option, wherein $\hat{\sigma}$ corresponds to the sample standard
deviation of the transformed training response values. The second is the “linear model” choice, where $\hat{\sigma}$ is based on the residual standard deviation from a least-squares linear regression
of scaled Y on the scaled X variables. Subsequently, a value
for $\nu$ within the range of 3–10 is selected to shape the prior
appropriately. Additionally, a suitable $\lambda$ value is determined such
that the qth quantile of the prior on $\sigma$ aligns with $\hat{\sigma}$ , that is
$P(\sigma < \hat{\sigma}) = q$. A suitable range for q is values between 0.7
and 0.99.}
\footnotetext[2]{Drawn independently for each dimension}
\footnotetext[3]{Given the prior centered at 0, representing our zero-centered transformed data}
\footnotetext[4]{These hyper-parameters are set such that $m\mu_\mu - k\sqrt{m}\sigma_\mu=-0.5$ and $m\mu_\mu + k\sqrt{m}\sigma_\mu=0.5$, where $k$ is a preselected hyperparameter determining the prior probability of $E[Y\mid\mathbf{x}, \mathbf{y}_\text{train}, \mathbf{X}_\text{train}]$}
\end{document}
